\documentclass[12pt,letterpaper]{article}
\usepackage[T1]{fontenc}
\usepackage{amsmath}
\usepackage{newtxtext,newtxmath}
\usepackage[letterpaper,margin=1in]{geometry}
\usepackage{longtable,booktabs,array}
\usepackage[hidelinks]{hyperref}

\newcommand{\F}{\mathbb{F}}
\newcommand{\Fp}{\mathbb{F}_p}
\newcommand{\X}{\#X}
\newcommand{\wt}{\operatorname{wt}}
\newcommand{\CNQ}{C_{NQ}}
\newcommand{\Cev}{C_{\mathrm{ev}}}

\title{Revision log for\\ ``On quadratic residue codes and hyperelliptic curves''\\[4pt]
\large Corrected and expanded version, Draft 1}
\author{David Joyner \and Claude (Anthropic)}
\date{November 11, 2026}

\begin{document}
\maketitle

\section*{Scope}
\noindent\textbf{Original:} \texttt{original/qrcodes-hyperell-dmtcs2.tex}, the
version submitted to DMTCS (received 20 October 2006), sole author D.~Joyner.\\
\textbf{Revised:} \texttt{qrcodes-hyperelliptic.tex}, D.~Joyner and Claude.
Current draft: Draft~1 (November 11, 2026).

\medskip
\noindent\textbf{Numbering.} Sections~\ref{sec:fmt}--\ref{sec:nov10} were written
before Draft~1 and use the numbering of that time. Draft~1 moved the literature
review to Section~2, so most numbers have changed. Section~\ref{sec:conc} gives a
concordance, and the map of results in Section~\ref{sec:map} uses the Draft~1
numbers.

\medskip
The published 2008 DMTCS version (doi:10.46298/dmtcs.429, arXiv:math/0609562v2) is a
\emph{different} revision of the same submission. This log compares only against
the 2006 file. In the original, numbering ran with one counter per environment; in
the revision, numbering is by section (for example, Theorem~4.2). The table in
Section~\ref{sec:map} maps old results to new ones.

Every numerical claim in either version was checked with the Python code in
\texttt{code/}, using exhaustive enumeration of $\CNQ$ for $p\le23$.

\section{Format and front matter}\label{sec:fmt}
\begin{enumerate}
\item Document class changed from \texttt{dmtcs} to \texttt{article}, 12pt, letter
paper, Times (\texttt{mathptmx}), compiled with \texttt{pdflatex}. The RCS macros,
\texttt{latin1} input encoding, \texttt{subfigure} and \texttt{natbib} were removed.
\item Author block: the affiliation ``Mathematics Dept., USNA, Annapolis, MD'' was
replaced by \texttt{wdjoyner@gmail.com}. Claude (Anthropic) was added as second
author, with a footnote and an acknowledgment paragraph saying what it contributed.
\item The abstract was rewritten to state the corrected results and constants.
Keywords were updated, and the MSC codes were updated to MSC 2020.
\item A table of contents was added.
\item The title-page date ``Original submission: 20 October 2006'' was replaced by
``November 2026''. The paragraph describing the relation to the 2006 version (at
first Section~1.5, \emph{About this revision}) was moved to an unnumbered section
just before the Acknowledgments. Apart from that paragraph, the revised paper no
longer refers to the 2006 version; all such comparisons are in this log.
\item The Acknowledgments now point to this revision log.
\item Layout: the weight-distribution table (Table~1) and Remark~4.7 were reset
so that the PDF has no overfull lines.
\item The material after \verb|\end{document}| (an unused remark with an unproved
lower bound on $|S_3|$) was deleted.
\end{enumerate}

\section{Mathematical corrections}\label{sec:math}

\subsection{Errors corrected}
\begin{enumerate}
\item \textbf{Lemma 1 (old), parity claims.} We have
$r_{S_1}^*r_{S_2}^*=r_{S_1}r_{S_2}+(|S_1|+|S_2|+1)J$, and
$\wt(r_{S_1}r_{S_2})\equiv|S_1||S_2|$. So in the cases (even, even) and (odd, odd)
the element $(r_{S_1}r_{S_2})^*$ has odd and even weight respectively, the opposite
of what was stated. The lemma was replaced by the short list of facts in \S2.1.
\item \textbf{``$*$ is an algebra involution on $R_{even}$.''} This is false: $*$ is
not additive, since $(a+b)^*=a+b+J$, and it maps even-weight elements to odd-weight
ones. Also, $R_{even}=(1+x)R$ contains no~$1$. The claim was deleted.
\item \textbf{Remark 1 (old).} ``$\CNQ$ contains a QR code'' was corrected to
``$\CNQ\subseteq\mathrm{QR}\oplus\mathrm{QR}'$'' (Remark~3.12).
\item \textbf{Lemma 2 (old)} relied on self-duality, which was stated only as a
conjecture (Conjecture 3, old). Self-duality is now proved (Theorem~3.4), so the
lemma is justified. The bound is now attributed correctly to Rains~\cite{Ra98}, not
Sloane--Mallows.
\item \textbf{Corollary 1 (old), $\max\X_S>\frac53p-4$.} The proof assumed that a
minimum-weight codeword comes from an even-size $S$. For $p=11$ it does not: the 77
weight-6 words all come from odd $S$. The result was restated and reproved as
Corollary~4.5, $M_p\ge\frac53p-6$ for $p\equiv3\pmod 4$ and $p\ge11$, using the new
twist identity.
\item \textbf{Remark after Proposition 1 (old).} The text gave both
``$\sqrt p+1$'' and ``$\sqrt p+2$'', and applied part~(c) to the wrong parity. The
corrected statement is Remark~3.12. That remark also notes that the bound follows
from Weil's bound and the trivial count. (The referee report said this bound was
stronger than Weil; that was wrong and is corrected here.)
\item \textbf{Theorem (old 2), proof.} The entropy function had a spurious leading
$\delta$. The weight formula now cites Proposition~1, not Lemma~3.
\item \textbf{Theorem (old 3), proof.} The proof claimed to contradict ``$B(1.62,p)$
for all sufficiently large $p$''. The correct logic is that $B(1.62,p)$ can hold for
only finitely many $p$. Rewritten as Theorem~4.3 with the sharp constant
$c^*=1.62584\ldots$.
\item \textbf{Theorem 4 (old, LQR parameters).} For $p\equiv3\pmod4$, the claimed
lower bound $d\ge4p-2\max\X_S$ was not proved. The injectivity argument was
incorrect: it asserted $r_Qr_{T^c}=0$, and it concluded ``$S_1=\emptyset$,
$S_2=\Fp$'', contradicting the injectivity it claimed. For $p\equiv1\pmod4$, $d_p$
should be $4p+4-2M$. The statement ``this code is non-linear'' is false for
$p\equiv1\pmod4$.
\item \textbf{Remark after Theorem 4 (old).} The remark doubted that the minimum
distance could be below $2p$. In fact it always is: the minimum distance is
$2d(\CNQ)$, for example $8<14$ at $p=7$.
\item \textbf{Lemma 7 (old).} For $p\equiv1\pmod4$ and odd $S$, the sum should be
over $f_{S^c}$.
\item \textbf{Example (old, LQR, $p=11$).} The two sets $S=\{1,2,3,4\}$ and
$\{1,2,3\}$ were interchanged. The section was removed (see Section~\ref{sec:removed}).
\item \textbf{Section 3 (old), genus.} The formula $\gamma=n+k+1-d$ was corrected to
$\gamma=n+1-k-d$.
\item \textbf{Example (old, $p=13$ zeta).} The text said ``only 8 of the 12 zeros''.
$P(T)$ has degree $16$, and exactly $12$ of its $16$ zeros lie on
$|T|=1/\sqrt2$. The coefficients were correct.
\item \textbf{Optimality of the $[26,13,6]$ code.} The claim rested on general
linear-code tables, which, as we understand them, allow $d=7$ for $[26,13]$. It was
restated as optimality among even codes, with an explicit caveat that this has not
been verified independently.
\item \textbf{Voloch's Theorem (old 6).} The constant window $-\frac12\le a\le\frac52$
holds only for $\ell=2$. It was replaced by $O(1)\in[-2,3]$, and ``$\ell>2$'' versus
``$\ell\ge2$'' was made consistent. The introduction's claim that $\X_Q>1.5p$ for
$p\equiv3\pmod8$ was corrected to $\X_Q=1.5p-\frac12$.
\item \textbf{Conjecture 3 (old).} ``$C\cap C^\perp=\emptyset$'' was corrected to
$\{0\}$, and ``direct product'' to ``direct sum''.
\end{enumerate}

\subsection{Former conjectures now proved}
\begin{enumerate}
\item Conjecture 3 (old): self-dual for $p\equiv3\pmod4$, LCD for $p\equiv1\pmod 4$.
Now Theorem~3.4, with a three-line proof via $\langle u,v\rangle=\mathrm{ct}(uv^-)$.
\item Conjecture 4 (old): $C'$ is formally self-dual with distribution $A+A^*$. Now
Theorem~6.2: $C'^\perp$ is the half-swap of $C'$, so $C'$ is isodual. The
distribution statement is immediate.
\item Section 3 (old) asked whether $\CNQ$ always satisfies RH for
$p\equiv3\pmod4$. Answered negatively: RH fails at $p=23$ (30 of 32 zeros on the
circle). This agrees with Boston--Hao~\cite{BH17}.
\end{enumerate}

\subsection{New results}
\begin{enumerate}
\item Twist identity (Lemma~2.4): $\X_T+\X_{T'}=2p+2$ when $\chi(f_T(a))=-1$, and
its consequence $M_p+m_p\ge2p+2$ (Corollary~2.6).
\item Even subcode $\Cev$ (Proposition~3.10). It has dimension $p-1$ for every $p$,
weights $2p+2-\X_T$, and $d(\Cev)=2p+2-M_p$. For $p\equiv3\pmod 4$ it is the
doubly-even subcode, so $\X_T\equiv0\pmod 4$ for even $T$.
\item Theorem~4.1 (Goppa): the congruence condition $p\equiv1\pmod4$ was removed,
and the threshold sharpened to $c<1.7799$. This supersedes the LQR result (old
Theorem 5, $c=1.57$).
\item Theorem~4.3 (MRRW): valid for all $p$, with $c^*=1.62584$.
\item Theorem~5.1(3): the case $p\equiv5\pmod 8$, $S=Q\cup\{0\}$, which gives
$\X_S=1.5p+\frac12$. Every prime $p\ge11$ now has an explicit $S$ with
$\X_S\ge1.5p-\frac12$.
\item New data: weight distributions for $p=17,19,23$ (Table~1), the RH status of
seven codes (Table~2), and $M_p$, $m_p$ for $p\le23$ (Table~4; Table~3 before
the addition of 10 November 2026).
\end{enumerate}

\section{Material removed}\label{sec:removed}
\begin{enumerate}
\item Section 4 (old), \emph{Long quadratic residue codes}: Lemmas 5--7, Theorem 4,
Theorem 5, and the LQR example. They were replaced by Remark~4.7, which explains
that the LQR code is a doubled code or a coset of one, and so has the same relative
distance as $\CNQ$ at half the rate. The even subcode gives the stronger result.
\item The ``Cyclotomic arithmetic'' section was condensed to the facts actually used.
The sets $H(S_1,S_2,a)$ and identity~(2) (old) were removed.
\item The quoted lemma ``Bazzi--Mitter Prop.~3.3'' (old Lemma 4) was absorbed into
Proposition~3.6.
\end{enumerate}

\section{Additions for readability}
\begin{enumerate}
\item Eleven worked examples (Examples 2.1, 2.3, 2.5, 3.3, 3.7, 3.8, 3.11, 6.1, 6.3
and two in Sections 4 and 5). They cover arithmetic in $R$, point counting via Legendre sums, the twist
identity, dimension, the weight formula, the even subcode, Voloch's construction
for $p=13$, and zeta polynomials.
\item Proofs were rewritten in full, including the dimension count
(Proposition~3.2) and Lemma~2.2 (Artin--Hasse--Weil) with the trivial
bounds~(3).
\item Section~8 (Section~7 before 10 November 2026), \emph{Open questions and
their difficulty}: a random-code
heuristic ($M_p\approx1.78p$) and an assessment of each open question.
\item Section~9 (formerly Section~8), \emph{Literature}, in two parts: before 2006
and 2006 and after.
\item Appendix~A, \emph{Glossary}: 54 entries with a short introduction. It is
placed after the Acknowledgments and before the bibliography.
\item An introductory paragraph was added at the start of every section and
subsection, saying what that part contains and what to expect.
\item In the list of results, ``answers Question~1.4 for'' now reads ``answers
Question~1.4 affirmatively for''.
\item Typos fixed, among others: McEliese $\to$ McEliece, Welsh $\to$ Welch,
Tsafsman $\to$ Tsfasman, Dordrechet $\to$ Dordrecht, ``This rest of'', ``the the'',
``if and only if and only if'', ``consquence'', ``we if''; the cross-reference
``Theorem \textbackslash ref\{cor:main\}''; and the inconsistent ``$p>5$''. The
second supplementary law is no longer called ``quadratic reciprocity''.
\end{enumerate}

\section{Bibliography}\label{sec:bib}
\begin{enumerate}
\item Updated: Bazzi--Mitter (IEEE IT 2006), Jiang--Vardy (IEEE IT 2004), Joyner
(DMTCS 2008).
\item Added: Alahmadi--\"Ozdemir--Sol\'e 2018, Boston--Hao 2017, Calderbank 1983,
Chen--Peterson--Weldon 1969, Chinen 2017/18, Grassl tables, Joyner--Kim 2011,
Karlin 1969, Kasami 1974, Mallows--Sloane 1973, Massey 1992, MRRW 1977, Rains 1998,
Shi--Xu--Sol\'e 2021.
\item Removed as uncited: SAGE 2.10, Wage 2006, Chinen's ``An abundance of invariant
polynomials''. The ``Wage'' remark on $|S|=2,3$ was dropped.
\item Bibliographic details marked from memory in the referee report should be
checked against the publishers' pages before submission.
\end{enumerate}

\section{Additions of 10 November 2026}\label{sec:nov10}
{\sloppy
These changes came out of a further discussion of the Bazzi--Mitter conjecture.
They add a section and some data; nothing earlier was withdrawn.
\begin{enumerate}
\item New Section~7, \emph{The split by $p$ modulo 8}. It is placed before
\emph{Open questions}, so the later sections are renumbered: Open questions is now
Section~8, Literature Section~9.
\begin{itemize}
\item Proposition~7.1: for $p\equiv\pm1\pmod8$, $\Cev$ is the direct sum of two
even-like binary QR codes, so $M_p=2p+2-d(\widehat{\mathcal Q}_p)$ exactly.
\item Proposition~7.2: for $p\equiv\pm3\pmod8$, $\Cev$ is $\F_4$-linear and is the
binary image of an even-like quaternary QR code, with $d_4\le d(\Cev)\le\frac43d_4$.
The bordered analog is due to Karlin, Bhargava and Tavares (1978).
\item Remark~7.3: the pure and bordered codes can differ ($p=13$: $d(\Cev)=8$,
$d(B_{13})=6$).
\item Theorem~7.4: Conjecture~1.3 implies linear minimum distance for binary or
quaternary QR codes along an infinite set of primes.
\item Conjecture~7.5 (new, refined): there is $c<2$ with $M_p\le cp$ for infinitely
many $p\equiv\pm3\pmod8$. Remark~7.6 records that it is equivalent to linear
distance for quaternary QR codes, that $c<1.7799$ would refute Goppa's conjecture,
and a stronger form for safe primes $p\equiv3\pmod8$.
\item Table~3 (new): exact $d(\Cev)$ and $M_p$ for all $p\equiv\pm3\pmod8$ with
$p\le83$, bounds for $p=107$, and exact $M_p$ for $p\equiv\pm1\pmod8$, $p<200$,
from the extended-QR distances of Tomlinson et al.\ (2017, Table~9.4).
\item Remark~7.7: the Legendre-symbol form $\X_T=|T|+2+2N_+(T)$, and the
outline of a possible proof that separates structured from unstructured sets.
\end{itemize}
\item Abstract: a closing sentence on the new section. List of results: item~8.
After Conjecture~1.3: a pointer to Section~7.
\item Open questions: the heuristic paragraph now says that the $1.78p$ heuristic
fits only $p\equiv\pm3\pmod8$. The Bazzi--Mitter paragraph cites Theorem~7.4. New
paragraphs on Conjecture~7.5 and on analytic hypotheses (GRH, Hooley).
\item Literature: Karlin--Bhargava--Tavares 1978 and Gaborit 2002 (before 2006);
Tomlinson--Tjhai--Ambroze--Ahmed--Jibril 2017 (after 2006). The About-this-revision
paragraph and the Acknowledgements mention the new section.
\item Glossary: new entries \emph{Binary image}, \emph{Bordered double-circulant
code}, \emph{Even-like code}, $N_\pm(T)$; the entry \emph{QR code} now covers
quaternary QR codes.
\item Bibliography: added Gaborit 2002, Hooley 1967, Karlin--Bhargava--Tavares 1978,
Tomlinson et al.\ 2017.
\item Code: \texttt{code/check\_mod8.py} checks Propositions~7.1 and~7.2 for small
$p$; \texttt{code/dmin\_cev.c} computes $d(\Cev)$ exactly for
$p\equiv\pm3\pmod8$ (Table~3).
\end{enumerate}\par}

\section{Draft 1: revision for the style guide (November 11, 2026)}\label{sec:draft1}

\subsection{Scope}
Claude brought the paper into line with D.~Joyner's style guide
(\texttt{joyner-math-paper-style-guide.md}, Draft~2 of November 11, 2026), following
his instructions of the same day. Those instructions forbade mathematical changes.
No statement, proof, hypothesis or status label was altered. All new mathematical
text was confined to what the guide requires, and it is listed in
Section~\ref{sec:newtext} so that it can be refereed. Semicolons in existing text
were left in place. Rule~3.8 (no semicolons) was applied to new text only. The
guide has no paper-specific rule (its Section~9) for this paper.

Before the changes, the paper compiled with no errors, no undefined references and
no overfull boxes. There were 7 hyperref warnings, all caused by math in the title
of the subsection on the extended code. The typewriter font was a bitmap font.
After the changes, the paper compiles with no errors and no warnings, and every font
is Type~1.

\subsection{Defaults used for the rules marked [DECIDE]}
\begin{enumerate}
\item Document class: the existing class, \texttt{article}, was kept.
\item Times package: \texttt{newtxtext} and \texttt{newtxmath}.
\item Glossary placement: notation near the beginning (a table at the end of
Section~3), and the glossary at the end, just before the Acknowledgments.
\item AI systems on the author line: the instructions did not mention the author
line, so it was left as it was (D.~Joyner and Claude). The Acknowledgments also
credit Claude and point to this log.
\end{enumerate}

\subsection{Changes}
\paragraph{Typesetting.}
\begin{enumerate}
\item \texttt{mathptmx} was replaced by \texttt{newtxtext,newtxmath}. The package
\texttt{amssymb} was dropped because it conflicts with \texttt{newtxmath}, which
supplies the same symbols. The same change was made in this log.
\item Equations are numbered within sections (\verb|\numberwithin|), as rule~1.6
requires. So (1)--(4) became (3.1), (3.2), (3.3) and (4.1).
\item \verb|\mathrm{circ}| became \verb|\operatorname{circ}| (rule~1.8).
\item The title of the subsection on the extended code uses
\verb|\texorpdfstring|, which removes the 7 hyperref warnings.
\item The hand-typed ``Item~6'' in Section~1 became a \verb|\ref| to a new label
\texttt{item:voloch} (rule~1.5).
\end{enumerate}

\paragraph{Title page, abstract and acknowledgments.}
\begin{enumerate}
\item The subtitle ``Revised and updated'' became ``Corrected and expanded
version. Draft 1'' (rule~6.2). The title itself is unchanged.
\item The date line ``November 2026'' became ``Original draft: 2006. Corrected and
expanded draft: 2026.'' (rule~6.2).
\item A \verb|\thanks| note on the title identifies the 2006 DMTCS submission and
the published revision~[Jo08], and gives the date of Draft~1 (rule~6.3).
\item The abstract's first sentence, ``This AI-aided paper revises and updates the
first author's earlier paper [Jo08]'', was replaced by an AI-disclosure sentence
that refers to the Acknowledgments (rule~5.2). The old sentence was also inaccurate:
the revision is of the 2006 submission, not of~[Jo08] (see Scope at the top of
this log). The rest of the abstract is unchanged.
\item ``Acknowledgements'' became ``Acknowledgments'' in the heading, in the
footnote on Claude, and in this log (rule~1.4).
\item The Acknowledgments now say that Claude acted as referee, reviser and
co-author (rule~5.3). They describe Claude's role in Draft~1, state that AI review is
not an independent certification of the results (rule~5.4), and say that further
details on the AI usage are in this log (rule~5.5). They give the commit hash of the
scripts (rule~8.4).
\end{enumerate}

\paragraph{Citations.} Every \verb|\bibitem| now prints its key as its label
(for example [Jo08] rather than [17]), and the bibliography is sorted by label
(rule~1.7). Five keys were renamed, and every \verb|\cite| was updated. The new keys
were chosen by D.~Joyner.
\begin{center}
\begin{tabular}{@{}lll@{}}
\toprule
Old key & New key & Note\\
\midrule
\texttt{Gr} & \texttt{Gr26} & ``accessed 2026'' added to the entry\\
\texttt{Vo} & \texttt{Vo04} & entry completed, see below\\
\texttt{TT17} & \texttt{TTAAJ17} & initials of all five authors\\
\texttt{Coy06} & \texttt{Co06} & two-letter initial, as for other single authors\\
\texttt{Ch18} & \texttt{Ch17} & matches the year of the arXiv entry cited\\
\bottomrule
\end{tabular}
\end{center}
The entry \texttt{Vo04} was a bare ``preprint''. It now reads: J.~F.~Voloch,
Asymptotics of the minimal distance of quadratic residue codes, in \emph{Finite
Fields: Theory and Applications}, Oberwolfach Report No.~54/2004, 2946--2947. These
details come from the copy of that report that D.~Joyner pointed to
(\url{https://ems.press/content/serial-article-files/45972}). The entry
\texttt{Ch17} still cites arXiv:1709.03389 (2017). That paper was later published in
\emph{Appl.\ Algebra Engrg.\ Comm.\ Comput.}\ 30 (2019), 275--284, and the entry
could be updated to the journal version.

\paragraph{Structure (Section 2 of the guide).}
\begin{enumerate}
\item The literature review (old Section~9) moved to Section~2, just after the
Introduction. Its text is unchanged. It was already split at 2006, the year of the
original manuscript.
\item The old Sections~2--8 became Sections~3--9.
\item Section~3 (\emph{Notation and background}) ends with a new table of notation
(Subsection~3.4).
\item The old Section~8, \emph{Open questions and their difficulty}, became
Section~9, \emph{Summary of results and open questions}. It begins with a new
orienting paragraph and a new Subsection~9.1, \emph{Summary of results}. The old
text forms Subsection~9.2, under the old title. Its first sentence now says ``This
subsection'' instead of ``This section''.
\item The glossary (Appendix~A) now comes before the Acknowledgments, not after
them.
\item The unnumbered section \emph{About this revision} was removed from the paper
(rule~6.4). Its content is now in the \verb|\thanks| note on the title and in this
log. Its text was:
\begin{quote}\small
This is a revision of the version submitted to \emph{Discrete Mathematics and
Theoretical Computer Science} on 20 October 2006. A different revision was
published as [Jo08] (see also arXiv:math/0609562). The present text corrects
several errors in the 2006 version, replaces some arguments, proves two statements
that were stated as conjectures, and adds examples, a literature survey and a
glossary. Section~7, with the refined Conjecture~7.5, was added on 10 November 2026.
A separate revision log lists every change.
\end{quote}
(In Draft~1 numbering, Section~7 and Conjecture~7.5 are Section~8 and
Conjecture~8.7.)
\end{enumerate}

\subsection{New text to be refereed}\label{sec:newtext}
Rules~3.2--3.5 required the following additions. All numbers are Draft~1 numbers.
Each numerical claim in N1--N4 is checked by
\texttt{code/check\_examples\_draft1.py} (commit \texttt{c6e9125}). No new result
is claimed. The examples illustrate results already in the paper.
\begin{enumerate}
\item[N1.] \textbf{Example~1.3} (\texttt{ex:Bp7}), after Definition~1.2 (rules~3.3
and~3.4). For $p=7$ it computes $\#X_N=12$ by hand. An exhaustive search shows that
$12$ is the largest $\#X_S$ over all $127$ non-empty $S$, attained also by the even
set $\{0,1,2,4\}$. So $B(c,7)$ and $B_{\mathrm{ev}}(c,7)$ hold exactly when
$c\ge12/7$.
\item[N2.] \textbf{Example~8.2} (\texttt{ex:pm1p7}), after Proposition~8.1 (rules~3.3
and~3.4). For $p=7$, $r_QR_{\mathrm{ev}}$ is $0$ and the seven cyclic shifts of
$x^2+x^4+x^5+x^6$, a $[7,3,4]$ code. So $d(\Cev)=4$ and $M_7=12$, and
$\widehat{\mathcal Q}_7$ is the $[8,4,4]$ extended Hamming code.
\item[N3.] \textbf{Example~8.4} (\texttt{ex:pm3p5}), after Proposition~8.3 (rules~3.3
and~3.4). For $p=5$, every nonzero word of $\Cev$ has $\F_4$-weight $4$ and binary
weight $4$ or $6$. So $d_4=4$, $d(\Cev)=4$ and $M_5=8$.
\item[N4.] \textbf{Example~9.1} (\texttt{ex:heuristic}), in Subsection~9.2 (rule~3.3).
It evaluates the heuristic constant $2(1-H^{-1}(1/2))=1.7799\ldots$ at $p=11$
and $p=23$ (predictions $19.6$ and $40.9$) against $M_{11}=16$ and $M_{23}=40$.
\item[N5.] \textbf{Orienting paragraphs} at the start of Subsections~8.1, 8.2 and~8.3
(rule~3.2). They summarize what follows. The one for 8.2 says that the binary
minimum distance of $\Cev$ is determined up to a factor of at most $4/3$, which
restates Proposition~8.3.
\item[N6.] \textbf{Summary of results} (Subsection~9.1) and the opening paragraph of
Section~9 (Section~2.7 of the guide). Seven items restate the main results with
references. Each states the status of its result as the body does. Item~4 marks
Theorem~5.1 as conditional, item~6 marks the zeta results as computational, and
item~7 keeps Conjecture~8.7 a conjecture.
\item[N7.] \textbf{Table of notation} (Subsection~3.4), with an orienting paragraph
(Section~2.5 of the guide). It has no caption, so the table numbers are unchanged.
\item[N8.] \textbf{Definitions at first use} (rule~3.5) in the list of results in
Subsection~1.3. Item~1 defines ``isodual'', item~3 defines $H^{-1}$, and item~8
defines $\Cev$, $M_p$ and $\widehat{\mathcal Q}_p$, which were used there before
their definitions.
\item[N9.] \textbf{Front matter}: the \verb|\thanks| note, the first sentence of
the abstract, and the new Acknowledgments paragraphs described above.
\end{enumerate}

\subsection{Concordance of numbers}\label{sec:conc}
The left column gives the numbering just before Draft~1, and the middle column gives
the Draft~1 numbering. Every \verb|\label| was kept. The table was generated from
the two source files and checked against the label numbers that \LaTeX\ wrote to
the \texttt{.aux} files. Sections map as follows: old~1 to new~1, old~2--8 to
new~3--9, and old~9 (Literature) to new~2. Table numbers did not change.
\begin{longtable}{@{}lll@{}}
\toprule
Before Draft 1 & Draft 1 & Label\\
\midrule\endhead
Conjecture 1.1 & Conjecture 1.1 & \texttt{conj:goppa}\\
Definition 1.2 & Definition 1.2 & \texttt{def:B}\\
Conjecture 1.3 & Conjecture 1.4 & \texttt{conj:BM}\\
Question 1.4 & Question 1.5 & \texttt{q:1}\\
Question 1.5 & Question 1.6 & \texttt{q:2}\\
(1) & (3.1) & \texttt{eq:minus}\\
(2) & (3.2) & \texttt{eq:wtprod}\\
Example 2.1 & Example 3.1 & \texttt{ex:p7arith}\\
Lemma 2.2 & Lemma 3.2 & \texttt{lem:AHW}\\
(3) & (3.3) & \texttt{eq:trivial}\\
Example 2.3 & Example 3.3 & \texttt{ex:p7points}\\
Lemma 2.4 & Lemma 3.4 & \texttt{lem:twist}\\
Example 2.5 & Example 3.5 & \texttt{ex:twist}\\
Corollary 2.6 & Corollary 3.6 & \texttt{cor:Mm}\\
Definition 3.1 & Definition 4.1 & \texttt{def:QQR}\\
Proposition 3.2 & Proposition 4.2 & \texttt{prop:dim}\\
Example 3.3 & Example 4.3 & \texttt{ex:dim}\\
Theorem 3.4 & Theorem 4.4 & \texttt{thm:duality}\\
(4) & (4.1) & \texttt{eq:ip}\\
Remark 3.5 & Remark 4.5 & (no label)\\
Proposition 3.6 & Proposition 4.6 & \texttt{prop:main}\\
Example 3.7 & Example 4.7 & \texttt{ex:p7weights}\\
Example 3.8 & Example 4.8 & \texttt{ex:p11}\\
Definition 3.9 & Definition 4.9 & (no label)\\
Proposition 3.10 & Proposition 4.10 & \texttt{prop:Cev}\\
Example 3.11 & Example 4.11 & \texttt{ex:Cev11}\\
Table 1 & Table 1 & \texttt{tab:wd}\\
Remark 3.12 & Remark 4.12 & \texttt{rem:sqrt}\\
Theorem 4.1 & Theorem 5.1 & \texttt{thm:goppa}\\
Remark 4.2 & Remark 5.2 & (no label)\\
Theorem 4.3 & Theorem 5.3 & \texttt{thm:mrrw}\\
Remark 4.4 & Remark 5.4 & (no label)\\
Corollary 4.5 & Corollary 5.5 & \texttt{cor:53}\\
Example 4.6 & Example 5.6 & (no label)\\
Remark 4.7 & Remark 5.7 & \texttt{rem:LQR}\\
Theorem 5.1 & Theorem 6.1 & \texttt{thm:voloch}\\
Example 5.2 & Example 6.2 & (no label)\\
Theorem 5.3 & Theorem 6.3 & \texttt{thm:chebotarev}\\
Corollary 5.4 & Corollary 6.4 & (no label)\\
Example 6.1 & Example 7.1 & \texttt{ex:zeta7}\\
Theorem 6.2 & Theorem 7.2 & \texttt{thm:isodual}\\
Example 6.3 & Example 7.3 & \texttt{ex:zeta13}\\
Table 2 & Table 2 & \texttt{tab:RH}\\
Proposition 7.1 & Proposition 8.1 & \texttt{prop:pm1}\\
Proposition 7.2 & Proposition 8.3 & \texttt{prop:pm3}\\
Remark 7.3 & Remark 8.5 & (no label)\\
Theorem 7.4 & Theorem 8.6 & \texttt{thm:mod8}\\
Conjecture 7.5 & Conjecture 8.7 & \texttt{conj:refined}\\
Remark 7.6 & Remark 8.8 & \texttt{rem:refined}\\
Table 3 & Table 3 & \texttt{tab:mod8}\\
Remark 7.7 & Remark 8.9 & \texttt{rem:legendre}\\
Table 4 & Table 4 & \texttt{tab:Mp}\\
--- & Example 1.3 & \texttt{ex:Bp7} (new)\\
--- & Example 8.2 & \texttt{ex:pm1p7} (new)\\
--- & Example 8.4 & \texttt{ex:pm3p5} (new)\\
--- & Example 9.1 & \texttt{ex:heuristic} (new)\\
\bottomrule
\end{longtable}

\subsection{Rules not applied, and why}
\begin{enumerate}
\item Rule~1.6, second sentence (do not number subsections that do not need a
number). All subsections are listed in the table of contents, and this log cites
them by number, so they stay numbered.
\item Rule~3.3 in Section~2 (Literature). A worked example there would add nothing.
\item Rule~3.8 in existing text, by instruction. Existing semicolons were kept.
\item Rules~3.6 and~3.7 (hypotheses before theorems, no skipped steps). Applying
them would change statements or proofs, which this draft did not do. No violation
was looked for.
\item Rules~4.3 and~4.4. No citation was checked against its source, and no
computation in the body was rerun, except as described in the next subsection.
\item Rule~6.6 (``Remark 2026.$n$'' for anything not fully verified). Applying it
would add status labels, which this draft did not change. One candidate is the
optimality of the $[26,13,6]$ code after Example~7.3, which the text already
qualifies in prose.
\item Rule~1.5 for citations typed by hand. ``Gleason--Prange'' (proof of
Proposition~8.1), ``Prange (1957--1959)'' and ``Gilbert (1952) and Varshamov
(1957)'' (Section~2) have no \verb|\cite|. Fixing them needs new bibliography
entries or precise theorem numbers that were not checked.
\item Rules~8.5 and~8.6 concern the repository, not the paper. The repository has
no license file. It is public and contains a corrected version of a manuscript
whose different revision was published as~[Jo08]. Neither point was acted on.
\end{enumerate}

\subsection{Not verified}
\begin{enumerate}
\item The Voloch abstract (\texttt{Vo04}) was read only through a web tool's
summary, not in a printed copy. That summary says the abstract shows that some
subfamilies of binary QR codes are asymptotically bad. Section~2 describes that
paper only as studying ``the asymptotics of QR code minimum distances'', and
Section~8 says that whether binary QR codes are asymptotically good is open. These
statements look compatible, but a referee should compare them with the abstract.
\item The new notation table and summary were checked against the body by reading,
not by an independent referee.
\item The bibliographic details of entries other than \texttt{Vo04} were not
rechecked in Draft~1.
\end{enumerate}

\section{AI usage}\label{sec:ai}
All AI work on this revision was done by Claude (Anthropic), recorded as Claude
Opus~5.5 in the paper.
\begin{enumerate}
\item \emph{Referee.} Claude refereed the 2006 manuscript. Its findings are the
corrections in Section~\ref{sec:fmt} onwards.
\item \emph{Co-author and reviser.} Claude proposed and wrote the proofs and
computations listed in the Acknowledgments of the paper, including the 10 November
2026 additions (Section~\ref{sec:nov10}). D.~Joyner reviewed them.
\item \emph{Programmer.} Claude wrote the scripts in \texttt{code/}, which check
every numerical claim.
\item \emph{Style reviser (Draft~1).} Claude made the changes in
Section~\ref{sec:draft1} and wrote the new text N1--N9, which has not yet been
refereed.
\end{enumerate}
AI review is not an independent certification of the results.

\section{Sources not consulted}\label{sec:notconsulted}
Rule~4.3 asks that every citation be checked against its source. In Draft~1 no
cited source was consulted, apart from the Oberwolfach report for \texttt{Vo04}
(through a web tool) and the publication record of \texttt{Ch17}. The earlier
referee report marked some bibliographic details as written from memory (see
Section~\ref{sec:bib}, item~4). The optimality claim for the $[26,13,6]$ code rests
on the tables~[Gr26] as the authors understand them, without independent
verification.

\section{Original and new results}
Rule~6.5 asks that original results be separated from results first established in
the revision. Section~\ref{sec:math} does this. Its subsections list the errors
corrected in original results, the former conjectures now proved, and the new
results. The map of results in Section~\ref{sec:map} gives the correspondence.
Draft~1 adds no new results.

\section{Map of results}\label{sec:map}
\begin{longtable}{@{}p{0.42\textwidth}p{0.52\textwidth}@{}}
\toprule
2006 version & Draft~1\\
\midrule\endhead
Lemma 1 (products) & \S3.1 facts (corrected)\\
Lemma 2 (mindist bound) & used in Corollary 5.5 (Rains)\\
Lemma 3 (Artin--Hasse--Weil) & Lemma 3.2\\
Lemma 4 (Bazzi--Mitter) & absorbed into Proposition 4.6\\
Proposition 1 & Proposition 4.6\\
Corollary 1 ($\frac53p-4$) & Corollary 5.5 ($\frac53p-6$, new proof)\\
Theorem 1 (Bazzi--Mitter) & Theorem 5.1 (no congruence condition)\\
Theorem 2 ($1.77$, $p\equiv1\bmod4$) & Theorem 5.1 ($c<1.7799$, all $p$)\\
Theorem 3 ($1.62$, $p\equiv1\bmod4$) & Theorem 5.3 ($c<1.62584$, all $p$)\\
Corollary 2 (1.62, all $p$) & Theorem 5.3\\
Conjecture 3 (self-dual / LCD) & Theorem 4.4 (proved)\\
Conjecture 4 ($C'$) & Theorem 7.2 (proved)\\
Section 4: Theorem 4, Theorem 5, Lemmas 5--7 & Remark 5.7 (removed)\\
Lemma 8 (Voloch) & Theorem 6.1 (with new case $p\equiv5\bmod8$)\\
Theorem 6 (Voloch, Chebotarev) & Theorem 6.3\\
--- & Lemma 3.4, Corollary 3.6, Proposition 4.10, Theorem 7.2 (new)\\
Conjecture (Bazzi--Mitter), restated & Conjecture 8.7 (refined, $p\equiv\pm3\bmod 8$)\\
--- & Propositions 8.1, 8.3, Theorem 8.6 (new, 10 November 2026)\\
\bottomrule
\end{longtable}

\begin{thebibliography}{BH17}
\bibitem[BH17]{BH17} N. Boston and J. Hao, The weight distribution of quasi-quadratic
residue codes, arXiv:1705.06413.
\bibitem[Ra98]{Ra98} E.~M. Rains, Shadow bounds for self-dual codes, \emph{IEEE Trans.\
Inform.\ Theory} 44 (1998), 134--139.
\end{thebibliography}
\end{document}
